\SourceLeaf{77}%
\[
\int z\mathrm{d}s-\alpha\alpha T+\alpha\alpha U
\]
i. e. eiusdem expressionis, quae minimum esse debet pro casu $\beta\beta=\alpha\alpha$. Hinc colligimus, figuram aequilibrii fluidi in vase, pro quo $\beta\beta>\alpha\alpha$, convenire cum figura aequilibrii eiusdem fluidi in vase, pro quo $\beta\beta=\alpha\alpha$, ea tamen differentia, ut illa in aequilibrio stricto desinere debeat in cutem crassitiei insensibilis. Ceterum ill. \textsc{Laplace} iam monuit, pro illo casu vas cute fluidi insensibilis crassitie obductum aequipollere vasi tali, cuius particulae vim attractivam in fluidum exerceant vi attractivae partium fluidi mutuae aequalem.

Sponte hinc sequitur modificatio, propositionibus art. 18 circa ascensum fluidorum in tubis capillaribus verticalibus adiicienda: quoties scilicet $\beta\beta>\alpha\alpha$, in formulis illic allatis $\alpha\alpha$ loco ipsius $\beta\beta$ substituere oportet.

\Article{33}
In casu eo, ubi $\beta\beta<\alpha\alpha$, madefactio vasis per cutem fluidi insensibilis crassitiei locum habere nequit, siquidem lex functionum $\theta',\Theta'$ ea est, ut valor functionis
\[
\alpha\alpha\left(1-\frac{\theta'\rho}{\theta'0}\right)-\beta\beta\left(1-\frac{\Theta'\rho}{\Theta'0}\right)
\]
pro qua brevitatis caussa scribemus $Q\rho$, continuo crescat, dum $\rho$ a valore $0$ versus valorem sensibilem progreditur: manifesto enim pro tali functionis $Q\rho$ indole existentia talis cutis conditioni minimi repugnaret. Sponte illam indolem affert hypothesis, de qua in art. 31 loquuti sumus, puta ubi $fx,Fx$ sunt in ratione determinata ab $x$ independente, quoniam hinc etiam sequitur $\frac{\theta'\rho}{\theta'0}=\frac{\Theta'\rho}{\Theta'0}$, et proin $Q\rho=(\alpha\alpha-\beta\beta)(1-\frac{\theta'\rho}{\theta'0})$. At si functiones $f,F$ legem diversam sequerentur, haud impossibile esset, ut valore ipsius $\frac{\theta'\rho}{\theta'0}$ rapidius decrescente, quam valore ipsius $\frac{\Theta'\rho}{\Theta'0}$, functio $Q\rho$. intra ambitum valorum insensibilium ipsius $\rho$, primo fieret negativa, et postquam attigisset valorem suum minimum (i. e. extremum negativum) rursus ascenderet per valorem $0$ versus limitem suum positivum $\alpha\alpha-\beta\beta$. In tali casu aequilibrium utique postularet cutem insensibilem, cuius crassities generaliter loquendo tanta esse deberet, ut $Q\rho$ haud sensibiliter discrepet a valore suo minimo. Qui si per $-\beta'\beta'$ denotatur, erit $\beta'\beta'<\beta\beta$; figura autem partis substantialis fluidi perinde determinabitur, ac si esset in vase, cuius respectu loco quantitatis $\beta\beta$ substituere oportet $\beta'\beta'$, i. e. angulus inter planum
\SourceLeaf{78}superficiem fluidi liberam in confiniis partis substantialis tangens atque parietem vasis erit $=2\operatorname{arc.}\sin\frac{\beta'}{\alpha'}$. Sed quum valde dubium sit, an talis casus in rerum natura exstet, superfluum videtur, diutius ei immorari.

\Article{34}
Alienum foret ab instituto nostro praesente, a principiis generalibus hic stabilitis ad phaenomena specialia descendere, praesertim quod illorum principiorum essentia quadrat cum theoria ea, per quam ill. \textsc{Laplace} aequali arte et successu permulta phaenomena in aequilibrio fluidorum conspicua iam explicavit. Vastus utique superest campus, largam messem novam pollicens: sed haec curis futuris reservata maneat. Contra e re erit, quasdam annotationes adiicere, quae vel novam lucem huic argumento affundere, vel interpretationem erroneam arcere poterunt.

I. Theoria nostra non arrogat sibi determinationem figurae aequilibrii mathematice exactam, sed acquiescit in determinatione figurae talis, a qua figura aequilibrii vera differre nequit quantitate sensibili. Errares, si hoc alicui imperfectioni theoriae tribueres, quae ex asse praestitit, quantum praestare possibile est, quamdiu lex attractionis molecularis ignoratur. In statu aequilibrii functio $\Omega$ exacte maximum esse debet, adeoque functio
\[
\frac{2\pi cs\psi0}g-\frac\Omega{gc}
\]
minimum; haec autem, pro indole attractionis molecularis, non quidem exacte aequalis est functioni $W$, attamen insensibiliter tantum ab ea differt. Figura igitur, pro qua $W$ fit minimum, non est exacta figura aequilibrii, sed differentia esse debet insensibilis, quatenus quidem quaelibet mutatio sensibilis istius figurae valorem sensibiliter maiorem functionis $W$ produceret. Manifesto hinc non excluditur differentia sensibilis in curvatura superficiei, dummodo limitetur ad partem superficiei insensibilem: quapropter in figura aequilibrii exacta angulum constantem supra per $A$ denotatum haud amplius considerare licet tamquam inclinationem superficiei fluidi ad parietem vasis in ipso contactu, sed tantummodo in distantia immensurabili a vase, sive, ut ill. \textsc{Laplace} recte iam monuit, inclinatio in limite sphaerae sensibilis attractionis vasis cum valore ipsius $A$ sensibiliter coincidet.

\SourceLeaf{79}II. Probe distinguere oportet figuram aequilibrii a figura quietis. Quoties fluidum est in statu aequilibrii, certo in eo perseverare debebit. At quoties figura fluidi aliquantum a figura aequilibrii differt, nihilominus accidere potest, ut fluidum vel in quiete permaneat, vel, si moveatur, motum iam amittat, antequam statum aequilibrii attigerit, perinde ut e. g. cubus plano horizontali tantum impositus in aequilibrio versatur, sed etiam supra planum inclinatum quiescere potest, frictione motum impediente. Ita fluidum talem statum occupans, pro quo $W$ habet valorem minimum, certo quiescet: sed quoties est in statu ab illo diverso, puta ubi $W$ diminutionis capax est, ex hoc statu in statum aequilibrii eatenus tantum transibit, quatenus frictio non impediverit. Hocce autem respectu duae conditiones aequilibrii essentialiter diversae sunt. Scilicet aequatio fundamentalis prior (art. 28) independens est a mutabilitate limitis $P$, i. e. ad conditionem minimi tunc quoque necessaria, ubi hic limes invariabilis supponitur: quapropter, quatenus quidem fluidum perfecta fluiditate gaudet, ut pars una supra alteram libere gliscere possit, dum vel minima vis motum postulat, fluidum necessario illi conditioni se accommodabit. Longe vero alia est ratio principii secundi (art. 29), quod essentialiter pendet a perfecta limitis $P$ mobilitate in superficie vasis. Conditio minimi in valore ipsius $W$ utique postulat aequationem $i=A$: si vero, postquam superficies fluidi priori quidem principio se accommodavit, angulus $i$ nondum assequutus est valorem normalem, neque adeo $W$ valorem absolute minimum, transitus in statum aequilibrii perfecti fieri nequit absque translocatione limitis $P$, sive absque motu fluidi in contactu cum vase, quali motui utique obstare potest frictio. Hinc manifestum est, cur in experimentis circa eadem corpora institutis tantas differentias in valore anguli $i$ offendamus. Perinde in casu eo, ubi $\beta\beta>\alpha\alpha$, fluidum in vase, cuius parietes iam sunt madefacti, utique se componet ad legem aequilibrii, secundum quam pro parte substantiali fluidi esse debet $i=180^0$: sed in vase, cuius parietes extra fluidum etiamnum sunt sicci, fluidum a statu non aequilibrato proficiscens parietesque vasis siccas invadens iam ad quietem pervenire poterit, antequam angulus $i$ valorem $180^0$ attigit. Hinc simul elucet ratio, cur phaenomena capillaria fluidorum talium, quae madefactioni non adversantur, in tubis siccis tantas irregularitates offerant, ascensumque saepissime longe minorem, quam in tubis iam humectatis, ubi consensum pulcherrimum cum theoria semper aspicimus.

\expandafter\def\csname HistoricalNoteText7\endcsname{Observare convenit, quantitatem ab ill. \textsc{Laplace} per $H$ denotatam convenire cum nostro $\pi c\theta0$, adeoque $a$ apud illum auctorem idem denotare, quod in signis nostris est $\frac{g}{\pi c\theta0}$ sive $\frac1{2\alpha\alpha}$.}%
\SourceLeaf{80}III. Ratio constantium $\alpha,\beta$ e phaenomenis derivari nequit, quoties $\beta$ est maior quam $\alpha$: figura enim eiusdem fluidi in vasibus forma aequalibus materia diversis pro isto casu non differt nisi respectu cutis immensurabilis vas madefacientis. Quoties autem $\beta$ minor est quam $\alpha$, determinatio rationis inter has constantes possibilis quidem est adiumento anguli $i$, sed propter rationes modo allatas magnam praecisionem vix feret. Pro mercurio in vasibus vitreis ill. \textsc{Laplace} statuit angulum $i=43^0\,12'$.

Longe maioris praecisionis capax est determinatio constantis $\alpha$, praesertim si vasibus madefactionem admittentibus uti licet. Pro aqua sub temperatura $8{,}5$ graduum thermometri centesimalis statuere oportet secundum experimenta ab ill. \textsc{Laplace} citata\HistoricalNote{7}
\[
\begin{aligned}
\alpha\alpha&=7{,}5675\ \text{millim. quadr., sive}\\
\alpha&=2{,}7509\ \text{millim.}
\end{aligned}
\]
Pro spiritu vini, cuius pondus specificum $=0{,}81961$, sub eadem temperatura
\[
\begin{aligned}
\alpha\alpha&=3{,}0441\ \text{millim. quadr., sive}\\
\alpha&=1{,}7447\ \text{millim.}
\end{aligned}
\]
Pro oleo terebinthino sub temperatura $8$ graduum
\[
\begin{aligned}
\alpha\alpha&=3{,}305\ \text{millim. quadr., sive}\\
\alpha&=1{,}818\ \text{millim.}
\end{aligned}
\]
Pro mercurio, sub temperatura $10$ graduum, statuere licet, donec experimenta nova maiorem praecisionem suppeditaverint,
\[
\begin{aligned}
\alpha\alpha&=3{,}25\ \text{millim. quadr., sive}\\
\alpha&=1{,}803\ \text{millim.}
\end{aligned}
\]
Ceterum verisimile est, temperaturam eatenus tantum valorem constantis $\alpha\alpha$ afficere, quatenus densitas inde pendet, cui itaque in hac hypothesi valor ipsius $\alpha\alpha$ proportionalis erit.

\SourceLeaf{81}Valores isti conclusi sunt ex ascensione vel depressione fluidorum in tubis capillaribus: attamen valde difficile est, horum diametros exacte mensurare, difficilius, de forma circulari sectionis transversalis certitudinem acquirere. Longe maiorem praecisionem pollicentur experimenta circa diametros et volumina magnarum guttarum mercurii tabulae horizontali vel curvaturae perparvae notae insidentium, qualia iam instituerunt physici \textsc{Segner} et \textsc{Gay-Lussac}: nec non, pro liquidis vasa vitrea madefacientibus, experimenta circa dimensiones bullarum magnarum aeris in vasibus superne operculo madefacto plano horizontali vel parum et secundum radium notum curvato clausis, ad quae instituenda physicos invitamus.

IV. Ne limites huic commentationi praescriptos excederemus, applicationem principiorum nostrorum generalium hocce quidem loco ad casum simplicissimum restringere oportuit, ubi liquidum unicum in vase firmo consideratur. Nihil vero obstat, quominus theoriae summa generalitas concilietur, ita ut etiam problema plurium liquidorum in eodem vase, nec non casum eum amplectatur, ubi insuper corpora rigida, vel omnino vel ex parte libera, fluido immersa sunt. Sed harum quaestionum uberiorem expositionem ad aliam occasionem nobis reservare debemus.
